% from ega4-17-fr.tex
\oldpage[IV-4]{93}

Ceci dit, l'application de (17.13.2) à la situation du diagramme (17.13.5.1) (compte tenu de (17.3.3, (iv)), qui implique que $X\times_SX$ est lisse sur $S$ en $x$ si $X$ l'est) donne le corollaire suivant.
\end{env}

\begin{corollary}[17.13.6]
\EGAReaderTarget{ega4.c04.cor.17-13-6}{IV.17.13.6}
Soient $S$ un préschéma, $X$, $Y$, $Z$ trois $S$-préschémas, $f:Y\to X$, $g:Z\to X$ deux $S$-morphismes; posons $T=Y\times_XZ$, et soient $p:T\to Y$, $q:T\to Z$ les projections canoniques. Soient $t$ un point de $T$, $y=p(t)$, $z=q(t)$, $x=f(y)=g(z)$. Supposons que $X$ soit lisse sur $S$ au point $x$, de dimension relative $m$, $Y$ (resp. $Z$) lisse sur $S$ au point $y$ (resp. $z$), de dimension relative $m+a$ (resp. $m+b$), $a$ et $b$ étant positifs ou négatifs. Alors les conditions suivantes sont équivalentes:
\begin{enumerate}
  \item[a)] $T$ est lisse sur $S$ au point $t$, de dimension relative $m+a+b$.
  \item[b)] L'homomorphisme
  \[
  \rho\otimes1:\Omega^1_{X/S}\otimes_{\sh{O}_X}k(t)\to
  (\Omega^1_{Y/S}\otimes_{\sh{O}_Y}k(t))\oplus
  (\Omega^1_{Z/S}\otimes_{\sh{O}_Z}k(t))
  \]
  où $\rho$ est donné par (17.13.5.4), est injectif.
  \item[c)] Le morphisme $T=Y\times_SZ\to X\times_SX$ est transversal à la diagonale de $X\times_SX$ au point $t$.
\end{enumerate}
Si $t$ est rationnel sur le corps résiduel de son image $s$ dans $S$, ces conditions sont aussi équivalentes à la suivante:
\begin{enumerate}
  \item[b')] L'homomorphisme
  \[
  \label{IV.17.13.6.1}
  \tag{17.13.6.1}
  \mathrm{T}_y(f)-\mathrm{T}_z(g):\mathrm{T}_{Y/S}(y)\oplus\mathrm{T}_{Z/S}(z)\to\mathrm{T}_{X/S}(x)
  \]
  (cf. (16.5.12.5)) est surjectif.
\end{enumerate}
De plus, lorsque les conditions a) et b) sont vérifiées en $t$, elles le sont dans un voisinage de $t$ dans $T$, et en restreignant $T$ à un tel voisinage, la suite
\[
\label{IV.17.13.6.2}
\tag{17.13.6.2}
0\to\Omega^1_{X/S}\otimes_{\sh{O}_X}\sh{O}_T\xrightarrow{\rho}
(\Omega^1_{Y/S}\otimes_{\sh{O}_Y}\sh{O}_T)\oplus
(\Omega^1_{Z/S}\otimes_{\sh{O}_Z}\sh{O}_T)
\xrightarrow{\sigma}\Omega^1_{T/S}\to0
\]
(où $\sigma$ est donné par (17.13.5.3)) est exacte.

Le seul point qui reste à prouver dans (17.13.6) concerne b'), et résulte de ce que, sous l'hypothèse que $t$ est rationnel sur $k(s)$, l'homomorphisme $\rho\otimes1$ est \emph{transposé} de $\mathrm{T}_y(f)-\mathrm{T}_z(g)$, en vertu de (17.13.5.4).

Lorsque les conditions équivalentes de (17.13.6) sont satisfaites, on dit que $f$ et $g$ forment un \emph{couple de morphismes transversaux au point $t$, relativement} à $S$; ici encore, on supprime souvent la mention de $S$.
\end{corollary}

\begin{env}[17.13.7]
\EGAReaderTarget{ega4.c04.statement.17-13-7}{IV.17.13.7}
Considérons en particulier le cas où $Y$ et $Z$ sont des \emph{sous-préschémas} de $X$, $f$ et $g$ étant les injections canoniques; $T$ est alors le sous-préschéma « intersection »

